\title{DeepSeekMath: Pushing the Limits of Mathematical Reasoning in Open Language Models}
Despite being trained on this limited dataset, DeepSeekMath-RL 7B surpasses DeepSeekMath-Instruct 7B in every evaluation metric, underscoring the significant impact of reinforcement learning. Table \ref{tab:base_math_pot_proof} demonstrates that DeepSeekMath-Base 7B attains robust results in proof autoformalization. We hypothesize that this limitation partly explains why our RL method primarily boosts Maj@K rather than other metrics. Notably, GRPO departs from the PPO method of inserting the KL penalty within the reward calculation; instead, it regularizes the objective by appending the KL divergence term between the candidate and reference policies, streamlining the calculation of $\hat{A}_{i,t}$. Through the development of a carefully engineered data selection process, we introduce the DeepSeekMath Corpus, a curated dataset comprising 120B tokens of mathematically relevant web content—approximately 7 times larger than the dataset used in Minerva \citep{minerva}, and about 9 times larger than OpenWebMath \citep{openwebmath}. Using only a selected portion of English instruction tuning data, GRPO delivers marked gains over the robust DeepSeekMath-Instruct across both in-domain (GSM8K: 82.9\% $\rightarrow$ 88.2\%, MATH: 46.8\% $\rightarrow$ 51.7\%) and out-of-domain (e.g., CMATH: 84.6\% $\rightarrow$ 88.8\%) mathematical evaluation sets during the reinforcement learning phase. The fastText model is trained with the following configuration: a vector dimension of 256, learning rate of 0.1, maximum word n-gram length of 3, a minimum word occurrence threshold of 3, and three epochs; this is implemented using the open-source toolkit\footnote{\url{https://fasttext.cc}}. Specifically, Table \ref{tab:code-math} illustrates that initializing training with code tokens produces a marked boost in the model’s performance on benchmarks like GSM8K and MATH with Python solutions; a further phase of math-specific training compounds these improvements. \item Our results show that GRPO significantly advances the capabilities of our instruction-tuned model DeepSeekMath-Instruct when using only instruction-tuning datasets. We propose that our methodological contributions to mathematical data curation offer a valuable foundation for further exploration in the research community, with ample opportunities for future advancements. Incorporating these new instances, the classifier undergoes retraining to achieve greater accuracy. Building on DeepSeek-Coder-Base-v1.5 7B \citep{deepseek-coder} as the initialization point, the model is trained for a total of 500B tokens. Moreover, combining code and math tokens in a joint, single-stage regime not only curtails the catastrophic forgetting encountered in sequential (two-stage) training but also strengthens results on both coding benchmarks (Table \ref{tab:code-math-general-eval}) and math reasoning tasks leveraging code (Table \ref{tab:code-math}). Additionally, we provide a cohesive framework for understanding multiple techniques and outline various avenues to optimize reinforcement learning approaches. The data is curated from Common Crawl (CC) leveraging a classifier based on fastText \citep{joulin2016fasttext}. These scores are normalized using the mean and standard deviation computed over the group. \subsection{Insights from Pre-Training}

We begin by discussing our pre-training experience. Additionally, HumanEval \citep{codex} and MBPP \citep{mbpp} are used to benchmark coding performance. \end{itemize}

\textbf{One-Stage Training}

\begin{itemize}[topsep=0pt]
    \item \textbf{Math Training for 150B Tokens}: The DeepSeek-LLM 1.3B model is directly trained on 150B tokens of math data;
    \item \textbf{Training on a mixture of 400B Code Tokens and 150B Math Tokens}: Considering that conducting math training after code training can lead to diminished code-related skills, we explore a single-stage approach by interleaving both code and math data during the same training run, assessing if this method can jointly enhance mathematical reasoning while preserving coding performance. Conversely, when code and math data are mixed for a single-stage approach, the model’s performance on pure mathematical reasoning tasks (without code) suffers. Although DeepSeekMath-Instruct matches the performance of leading proprietary Chinese models like GLM-4 and Baichuan-3 on MATH, it does not reach the level of GPT-4 or Gemini Ultra. According to the framework encapsulated in Equation \ref{eq:objective}, three essential factors—Data Source, Algorithm, and Reward Function—can be distinguished. 2) Notably, the training of DeepSeekMath-RL 7B was confined to chain-of-thought-format instruction tuning on GSM8K and MATH, using DeepSeekMath-Instruct 7B as the initialization. \textbf{Math Pre-Training at Scale}

\begin{itemize}[topsep=0pt]
    \item We demonstrate that the publicly available Common Crawl data possesses substantial value for mathematical applications. The advantage estimate, $A_t$, is obtained using Generalized Advantage Estimation (GAE) \citep{gae}, computed from rewards $\{r_{\ge t}\}$ and a parameterized value function $V_{\psi}$. As specified by Equations \ref{eq:GC-RFT} and \ref{eq:GC-GRPO}, a fundamental distinction exists between GRPO and Online RFT: GRPO dynamically modifies its gradient coefficient as a function of the reward values from the reward model, thereby enabling finer-grained positive or negative updates according to response quality. Each of these responses receives a score from the reward model, resulting in a corresponding set of reward values $\mathbf{r} = \{r_1, r_2, \cdots, r_G\}$. To probe the mechanisms behind RL’s effectiveness, we assess the Pass@K and Maj@K accuracies of both Instruct and RL models on two separate benchmarks. Remarkably, DeepSeekMath-Base leads on Chinese evaluations as well, likely due to our deviation from the strategy in prior studies \citep{minerva,llemma} that utilize exclusively English math data, and instead incorporate extensive high-quality datasets in other languages. \item We report on math-focused training experiments, establishing that pre-training on code data before mathematics leads to enhanced problem-solving abilities in mathematical tasks, regardless of tool use. On the challenging MATH dataset, it outperforms all existing open-source models and surpasses most proprietary competitors (such as Inflection-2 and Gemini Pro) by a minimum margin of 9 percentage points. Notably, mathematical pre-training enhances not only the model’s proficiency in mathematics but also bolsters performance on general reasoning tasks, as reflected in improvements across MMLU \citep{mmlu} and BBH \citep{bbh} benchmarks. Table \ref{tab:base_math_cot} presents results indicating that DeepSeekMath-Base 7B achieves top performance across all eight evaluation benchmarks when compared to open-source base models—including prominent models like Mistral 7B \citep{mistral} and the recently launched Llemma 34B \citep{llemma}, which has received mathematical training on Proof-Pile-2 \citep{llemma}. Future efforts will be directed towards refining our data selection strategy, aiming to curate even higher-quality pre-training corpora. \paragraph{Data Source} The data source underpins every training methodology. Beyond tool-assisted reasoning, code-focused pretraining also benefits direct mathematical problem solving. Additionally, we introduce a unified perspective for interpreting representative training methodologies, framing them as variants or simplifications of RL strategies. In this work, our approach relies exclusively on the question set from instruction tuning and applies straightforward nucleus sampling to generate responses. Subsequently, the reference model is reset to the current policy model, and training of the policy model proceeds using the updated reward model. \subsection{Validating the Quality of the DeepSeekMath~Corpus}

To evaluate the DeepSeekMath~Corpus, we perform pre-training experiments and compare its performance with recently available mathematics-focused corpora:

\begin{itemize}[topsep=0pt]
    \item \textbf{MathPile} \citep{mathpile}: This multi-source dataset, comprising 8.9B tokens, integrates content from textbooks, Wikipedia, ProofWiki, CommonCrawl, StackExchange, and arXiv, with over 85\% of the material coming from arXiv;
    \item \textbf{OpenWebMath} \citep{openwebmath}: Built by filtering CommonCrawl for mathematical data, yielding a total of 13.6B tokens;
    \item \textbf{Proof-Pile-2} \citep{llemma}: This corpus aggregates OpenWebMath, AlgebraicStack (10.3B tokens of mathematical code), and arXiv papers (28.0B tokens). Despite the strong quantitative reasoning abilities exhibited by DeepSeekMath, its proficiency in domains such as geometry and formal theorem-proving remains somewhat limited compared to leading closed-source systems. To assess the mathematical reasoning capabilities of the models, we conduct evaluations both in tool-free and tool-assisted scenarios, utilizing four quantitative reasoning benchmarks in both English and Chinese. Uncollected web pages connected to these URLs are subsequently appended to the seed corpus, supplying a richer set of positive examples and enabling the training of a more robust fastText model. This observation spans datasets for quantitative reasoning like GSM8K and MATH (see Table \ref{tab:arxiv-cot}), multiple-choice assessments such as MMLU-STEM (Table \ref{tab:arxiv-cot}), as well as tasks in formal mathematics, e.g., miniF2F (Table \ref{tab:arxiv_atp}). \section{Reinforcement Learning}

\subsection{Group Relative Policy Optimization}
Reinforcement learning (RL) has shown notable success in further advancing LLMs’ mathematical reasoning following Supervised Fine-Tuning (SFT) \citep{wang2023math,wizardmath}. \paragraph{Observation about Data Source} We categorize the origin of data into two primary types: online sampling and offline sampling. To rigorously evaluate the mathematical proficiency of DeepSeekMath-Base 7B, we analyze its capacity to generate fully self-contained solutions without auxiliary resources, tackle mathematical tasks requiring tool assistance, and engage in formal theorem proving. Pre-training on mathematical data contributes to enhancements in both understanding and reasoning abilities. In our experiments involving Proof-Pile-2, we maintain the arXiv:Web:Code token distribution ratio of 2:4:1 as outlined in \cite{llemma}. This offers evidence in support of the hypothesis that code pre-training contributes positively to reasoning skills, at least within the context of mathematical problem solving. Following instruction-based tuning and reinforcement learning, DeepSeekMath-Instruct and DeepSeekMath-RL attain strong results, exceeding 50\% accuracy on the challenging, competition-tier MATH benchmark—a milestone not previously reached by open-source models. In general, the gradient with respect to the parameters $\theta$ for any of these methods can be formalized as:

\begin{equation}
    \nabla_{\theta}\mathcal{J}_{\textcolor{red}{\mathcal{A}}}(\theta) = \mathbb{E}[\underbrace{(q,o) \sim \textcolor{red}{\mathcal{D}}}_{Data \ Source}]\left( \frac{1}{|o|} \sum_{t=1}^{|o|}  \underbrace{GC_{{\mathcal{A}}}(q, o, t, \textcolor{red}{\pi_{{rf}}})}_{Gradient \ Coefficient}  \nabla_{\theta}\log \pi_{\theta}(o_t | q, o_{<t})\right). Of particular note, DeepSeekMath-Base attains over a 10\% absolute performance gain on the competition-grade MATH dataset compared to all existing open-source base models, and it even surpasses the closed-source Minerva 540B \citep{minerva}—which, despite its much larger scale (77 times bigger), was pre-trained on PaLM \citep{palm} and further specialized with mathematical materials. Table \ref{tab:corpora_comparison} demonstrates that pretraining on DeepSeekMath~Corpus leads to increased mathematical reasoning proficiency in both English and Chinese. Consistent with the methodology of DeepSeek LLMs, optimization is carried out using AdamW \citep{adamW}, adopting $\beta_1=0.9$, $\beta_2=0.95$, and a $\mathrm{weight\_decay}$ of 0.1. GRPO eliminates the need for a separate critic model by estimating baselines directly from grouped scores, thus significantly conserving computational resources during training. PPO improves LLMs by maximizing the surrogate objective defined as:

\begin{equation}
\footnotesize
            \mathcal{J}_{PPO}(\theta) = \mathbb{E}{[q \sim P(Q), o \sim \pi_{\theta_{old}}(O|q)]} \frac{1}{|o|} \sum_{t=1}^{|o|} \min \left[ \frac{\pi_\theta(o_{t} | q, o_{<t})}{\pi_{\theta_{old}}(o_{t} | q, o_{<t})} A_{t}, \text{clip} \left( \frac{\pi_\theta(o_{t} | q, o_{<t})}{\pi_{\theta_{old}}(o_{t} | q, o_{<t})}, 1 - \varepsilon, 1 + \varepsilon \right)  A_{t} \right] ,
\end{equation}

Here, $\pi_{\theta}$ and $\pi_{\theta_{old}}$ denote the current and previous policy models, respectively, while $q$ and $o$ represent questions and responses sampled from both the question dataset and the prior policy $\pi_{\theta_{old}}$. \end{itemize}

\subsection{Summary of Evaluations and Metrics}

\begin{itemize}[topsep=0pt]
    \item \textbf{English and Chinese Mathematical Reasoning}: Our evaluation spans comprehensive testing on both English and Chinese mathematical benchmarks, ranging from elementary to collegiate-level problems. Each task requires models to write Python code for solution, with access to computational libraries like \textit{math} and \textit{sympy} for complex calculations. In total, our training dataset contains 776K instances. Illustrated in Figure \ref{fig:data_collect}, we implement an iterative process that systematically compiles a large-scale math dataset from Common Crawl, beginning with a high-quality, focused seed set (such as a small, curated mathematical data collection). Analysis of Figure \ref{fig:rl-analysis} reveals that Online RFT achieves substantial performance improvements over RFT across both evaluation benchmarks. The following representative approaches can be examined within this common framework:

\begin{itemize}
    \item \textbf{Supervised Fine-tuning (SFT)}: SFT adapts the pretrained model using manually curated SFT datasets. The policy is then updated to maximize the group-relative policy objective given in Equation~(\ref{eq:GRPO-obj}). The peak learning rate is set to 5.3e-4. Solutions for these problems are annotated in both CoT and tool-integrated reasoning forms. \paragraph{Natural Language Understanding, Reasoning, and Code}

We measure the model’s proficiency in natural language understanding using MMLU \citep{mmlu}, reasoning using BBH \citep{bbh}, and coding via the HumanEval \citep{codex} and MBPP \citep{mbpp} benchmarks. \item \textbf{Online Rejection Sampling Fine-tuning (Online RFT)}: Distinct from standard RFT, Online RFT commences with the SFT policy and refines it by training on outputs sampled dynamically from the current policy. In comparison, the majority of existing mathematical corpora are overwhelmingly English-focused, and tend to offer negligible or even detrimental effects on Chinese language mathematical reasoning. \item \textbf{Large-scale}:  
    In terms of scale, DeepSeekMath~Corpus dwarfs previously available mathematical datasets. This, in turn, facilitates improved retrieval of mathematical content in further collection cycles. \section{Supervised Fine-Tuning}

\subsection{SFT Data Curation}

We assemble a dataset for mathematical instruction tuning that encompasses both English and Chinese mathematical problems, sourced from various domains and exhibiting a wide range of difficulty. \label{eq:objective}
\end{equation}

This formulation consists of three fundamental elements: 1) \textit{Data Source $\mathcal{D}$}, which specifies the dataset utilized for training; 2) \textit{Reward Function $\pi_{{rf}}$}, which provides the reward feedback for model updates; 3) \textit{Algorithm $\mathcal{A}$}, which uses both the data and the reward signals to generate the gradient coefficient $GC$, regulating the degree of penalization or reward assigned during optimization. We further articulate a unified conceptual framework encompassing approaches such as Rejection Sampling Fine-Tuning (RFT) \citep{yuan2023scaling}, Direct Preference Optimization (DPO) \citep{dpo}, PPO, and GRPO. Moreover, for the KL term—unlike the implementation in (\ref{eq:PPO-reward})—we

\begin{algorithm}[t]
  \small
  \caption{Iterative Group Relative Policy Optimization}
  \textbf{Input} initial policy model $\pi_{\theta_{\text{init}}}$; reward models $r_\varphi$; task prompts $\mathcal{D}$; 
  hyperparameters $\varepsilon$, $\beta$, $\mu$
  \begin{algorithmic}[1]
    \State Initialize policy model: $\pi_\theta \leftarrow \pi_{\theta_{\text{init}}}$
    \For{iteration = 1, \dots, I}
       \State Set the reference model $\pi_{ref} \leftarrow \pi_{\theta}$
      \For{step = 1, \dots, M}
      \State Draw a mini-batch $\mathcal{D}_b$ from $\mathcal{D}$
      \State Update previous policy: $\pi_{\theta_{old}} \leftarrow \pi_{\theta}$ 
      \State For each prompt $q \in \mathcal{D}_b$, sample $G$ outputs $\{o_i\}_{i=1}^G$ from $\pi_{\theta_{old}} (\cdot \mid q)$
      \State Evaluate each generated output $o_i$ with the reward model $r_{\varphi}$ to obtain rewards $\{r_i\}_{i=1}^G$
      \State Calculate group-relative advantage estimates $\hat{A}_{i,t}$ for each token $t$ in $o_i$
      \For{GRPO iteration = 1, \dots, $\mu$}
        \State Update $\pi_{\theta}$ by maximizing the GRPO objective (see Equation \ref{eq:GC-GRPO})
      \EndFor
    \EndFor
    \State Continuously update $r_\varphi$ with replay-based training
    \EndFor 
  \end{algorithmic}
  \textbf{Output} $\pi_\theta$
  \label{alg:iter-grpo}
\end{algorithm}

\subsubsection{Outcome Supervision RL with GRPO} 
Consider the case where, given a prompt $q$, a set of responses $\{o_1, o_2, \cdots, o_G\}$ is drawn from the previous policy model $\pi_{\theta_{old}}$. The results shown in Table \ref{tab:understanding_reasoning_code} highlight marked improvements of DeepSeekMath-Base 7B over the previous DeepSeek-Coder-Base-v1.5 \citep{deepseek-coder} on MMLU and BBH, reflecting the beneficial effect of targeted mathematical training for language comprehension and reasoning tasks. Nonetheless, this finding should be interpreted cautiously given the scope of our experiments. The full Common Crawl dataset is segmented into separate domains, with each domain comprising web pages sharing an identical base URL. Unless specified otherwise, the training protocols follow those in Section \ref{sec:quality-policy}. Construction of reward model training data follows the methodology outlined in \citep{wang2023math}. We assess, for each domain, the proportion of web pages captured during the initial collection. Filtering of inferior mathematical web content is accomplished by ranking candidates according to their fastText-assigned scores and retaining only the top-scoring entries. Concretely, for each $q$, GRPO draws a collection of outputs $\{o_1, o_2, \cdots, o_G\}$ from the previous policy $\pi_{\theta_{old}}$, and then updates the policy parameters by maximizing the objective below:

\begin{equation}
\footnotesize
\begin{split}
    \mathcal{J}_{GRPO}(\theta) &= \mathbb{E}{[q \sim P(Q), \{o_i\}_{i=1}^G \sim \pi_{\theta_{old}}(O|q)]}  \\
    & \frac{1}{G}\sum_{i=1}^G\frac{1}{|o_i|} \sum_{t=1}^{|o_i|} \left\{ \min \left[ \frac{\pi_\theta(o_{i,t} | q, o_{i,<t})}{\pi_{\theta_{old}}(o_{i,t} | q, o_{i,<t})} \hat{A}_{i,t}, \text{clip} \left( \frac{\pi_\theta(o_{i,t} | q, o_{i,<t})}{\pi_{\theta_{old}}(o_{i,t} | q, o_{i,<t})}, 1 - \varepsilon, 1 + \varepsilon \right)  \hat{A}_{i,t} \right] - \beta \mathbb{D}_{KL}\left[\pi_{\theta} || \pi_{ref}\right]\right\} ,
\end{split}
\label{eq:GRPO-obj}
\end{equation}

In this context, $\varepsilon$ and $\beta$ are still hyper-parameters, and $\hat{A}_{i,t}$ is a group-wise advantage, which is derived using only the relative rewards within the group; we elaborate on this computation in later subsections. \item Our DeepSeekMath-Base 7B base model, after pre-training, matches the performance of Minerva 540B \citep{minerva}. Evaluation shows that our corpus maintains high standards, as evidenced by the DeepSeekMath-Base 7B model, which obtains 64.2\% on GSM8K \citep{gsm8k} and 36.2\% on the competition-level MATH benchmark \citep{MATH}, outperforming Minerva 540B \citep{minerva}. As a result, the PPO approach necessitates the concurrent training of a value function along with the policy. \paragraph{Mathematical Problem Solving with Tool Use} We assess model capability in program-supported mathematical problem solving on both GSM8K and MATH datasets via few-shot program-of-thought prompting \citep{pot,pal}. In GRPO, the learning rate for the policy model is 1e-6, with a KL regularization coefficient set to 0.04. \end{itemize}

Our experiments involved separate pre-training runs: DeepSeek-LLM 1.3B was trained on 150B tokens and DeepSeek-Coder-Base-v1.5 7B on 40B tokens for each respective arXiv dataset. We investigate this by training models at varying scales, including DeepSeek-LLM 1.3B and DeepSeek-Coder-Base-v1.5 7B \citep{deepseek-coder}, on arXiv datasets processed with different methodologies:

\begin{itemize}[topsep=0pt]
    \item \textbf{MathPile} \citep{mathpile}: an 8.9B-token dataset curated using cleaning and filtering heuristics, consisting of more than 85\% scientific arXiv publications;
    \item \textbf{ArXiv-RedPajama} \citep{redpajama}: the full collection of arXiv LaTeX documents stripped of preambles, comments, macros, and bibliographies, yielding 28.0B tokens in total. Through this framework, we categorize these techniques as variants of direct or streamlined reinforcement learning. Among the discussed algorithms, RFT and DPO employ an offline approach, whereas Online RFT and GRPO rely on online sampling. In the initial stage, the classifier is trained using samples from OpenWebMath \citep{openwebmath} as the positive set, complemented by a varied selection of unrelated web pages as negative instances. \end{itemize}

\subsection{Training and Evaluating DeepSeekMath-Base 7B}

This section details the development and evaluation of DeepSeekMath-Base 7B, a foundational model designed with robust mathematical reasoning skills. Together, these benchmarks span a variety of mathematical subfields, from elementary concepts up to those encountered at the undergraduate level. The initial reward model is fine-tuned from DeepSeekMath-Base 7B, using a learning rate of 2e-5. All training is implemented with the resource-efficient HAI-LLM framework \citep{haillm}. We further propose Group Relative Policy Optimization (GRPO), an adaptation of Proximal Policy Optimization (PPO), capable of substantially elevating mathematical reasoning while being more memory-efficient. Furthermore, since the DeepSeekMath Corpus encompasses multiple languages, we observe notable gains on Chinese mathematical benchmarks \citep{wei2023cmath,agieval}. Using OpenWebMath, we train a fastText classifier \citep{joulin2016fasttext} aimed at retrieving additional mathematical web documents resembling the seed set. With the exception of setting the maximum learning rate to 4.2e-4 and utilizing a batch size of 10M tokens, we largely adhere to the training protocol outlined in Section \ref{sec:quality-policy}. \item \textbf{Open-source models} span general-purpose systems—such as (1) DeepSeek-LLM-Chat 67B \citep{deepseek-llm}, (2) Qwen 72B \citep{qwen}, (3) SeaLLM-v2 7B \citep{seallm}, and (4) ChatGLM3 6B \citep{chatglm3}—as well as a number of models specialized for mathematical problem solving: (5) InternLM2-Math 20B~\footnote{\url{https://github.com/InternLM/InternLM-Math}}, which enhances InternLM2 via dedicated math training and subsequent instruction tuning; (6) Math-Shepherd-Mistral 7B, which utilizes PPO \citep{schulman2017proximal} on Mistral 7B \citep{mistral} guided by a process-supervised reward model; (7) the WizardMath suite \citep{wizardmath}, which applies evolve-instruct (AI-evolved instruction tuning) and PPO to boost mathematical reasoning in Mistral 7B and Llama-2 70B \citep{llama2} using problems chiefly from GSM8K and MATH; (8) MetaMath 70B \citep{metamath}, built on Llama-2 70B with further fine-tuning on an expanded GSM8K and MATH corpus; (9) ToRA 34B \cite{tora}, developed by adapting CodeLlama 34B for mathematical tool-augmented reasoning; and (10) MAmmoTH 70B \citep{MathInstruct}, a Llama-2 70B variant fine-tuned via instruction on MathInstruct. Our work partially addresses this question in the context of mathematics, finding that training with code improves models' mathematical reasoning capabilities, irrespective of tool usage. Within the two-stage approach, pretraining with code offers modest improvements that are compounded in subsequent math-focused training, yielding superior final performance. The model is trained over 500 iterations, with a batch size set at 256 and a fixed learning rate of 5e-5. The advantage for each token is then defined as the sum over normalized rewards from all subsequent reasoning steps: $\hat{A}_{i, t} = \sum_{index(j) \ge t} \widetilde{r}_i^{index(j)}$. We posit that adopting \textbf{WEAK-TO-STRONG} alignment frameworks \citep{burns2023weak} could drive transformative improvements in how learning algorithms operate. \end{itemize}

\item \textbf{Formal Mathematics}: To assess DeepSeekMath-Base on formalization tasks, we use the informal-to-formal theorem proving benchmark described in \citep{dsp_proof}, applied to miniF2F \citep{minif2f}, with Isabelle \citep{isabelle} selected as the proof assistant. \subsubsection{Code Training Enhances Mathematical Reasoning}

There exists a widely-held, yet largely untested, assumption that code-focused training enhances reasoning skills. \end{itemize}

These models are intended for a broad range of tasks, with most having completed various alignment processes. Moreover, the GRPO+PS variant exhibits better results than GRPO+OS, which points to the advantage of implementing fine-grained, step-specific gradient coefficients. During training, examples are concatenated in a random manner up to a context window of 4K tokens. \item \textbf{Chinese mathematical datasets}: We gather K-12 mathematics questions in Chinese that are distributed across 76 specific sub-domains, including topics like linear equations. For instance, even the PRM800K datasets \citep{lightman2023let}, meticulously annotated by skilled professionals, reportedly include nearly 20\% misannotations\footnote{\url{https://github.com/openai/prm800k/issues/12\#issuecomment-1728491852}}. Next, we manually curate the URLs representing mathematical content from these high-yield domains (for example, \url{mathoverflow.net/questions}). \paragraph{Observation about Gradient Coefficient} The reward signal is utilized by the algorithm to determine the gradient coefficient for updating model parameters. Motivated by \cite{wang2023math}, we incorporate process supervision, which offers feedback at each reasoning step. \paragraph{Reward Function} The reward function constitutes the primary mechanism by which training feedback is delivered. By evaluating self-consistency over 64 sampled outputs, DeepSeekMath~7B attains 60.9\% on MATH. This enhanced capacity for mathematical problem-solving in DeepSeekMath~can be ascribed to two primary factors: firstly, an advanced and rigorous data filtering pipeline that fully exploits the richness of open web-sourced datasets; secondly, the proposal of Group Relative Policy Optimization (GRPO)—a new adaptation of Proximal Policy Optimization (PPO)—which not only augments mathematical reasoning skills but also improves the memory efficiency of PPO. In large language model scenarios, only the final token typically receives a reward signal, which introduces difficulties in training an accurate value function at every token position. Each mathematical dataset is used to train a separate model for 150B tokens. The maximum context length is capped at 1024 tokens, and each training batch contains 1024 examples. \end{abstract}

\section{Introduction}

Recent progress in large language models (LLM) has fundamentally reshaped AI-driven mathematical reasoning, propelling significant strides on benchmarks designed to evaluate quantitative reasoning~\citep{MATH} as well as those focused on geometry~\citep{trinh2024solving}. \end{itemize}

Consequently, additional investigation is warranted, and we defer a comprehensive analysis to subsequent studies. To overcome these challenges, as visualized in Figure \ref{fig:grpo}, we introduce Group Relative Policy Optimization (GRPO), an approach that eliminates the necessity for an extra value function by leveraging the mean reward across a set of outputs, all sampled from the same question, as a baseline. Specifically, we designate 500,000 samples from OpenWebMath as positive examples, and randomly select another 500,000 Common Crawl web pages as negatives for the training. As described in Algorithm \ref{alg:iter-grpo}, iterative GRPO involves regenerating training datasets for the reward model using newly sampled outputs from the evolving policy model, while continually updating the reward model via a replay buffer that retains 10\% of earlier data. In reinforcement learning, this typically takes the form of a neural reward model. As illustrated in Figure \ref{fig:rl-analysis}, GRPO outperforms online RFT, underscoring the effectiveness of adjusting both positive and negative gradient coefficients. 
\begin{abstract}
Mathematical reasoning remains a formidable obstacle for language models, given its inherent complexity and highly structured requirements. This result not only exceeds the performance of all open-source models within the 7B to 70B parameter range but also outperforms most closed-source alternatives. With outcome supervision, the normalized reward for each output $o_i$ is provided only at its terminal state. This phenomenon can be attributed to the high similarity between the actor and SFT models early in training, which yields data with only marginal differences. By contrast, Online RFT treats all correct responses equally, reinforcing them with a uniform coefficient, and does not explicitly penalize incorrect outputs. The alternative corpora, by contrast, are limited in size, necessitating multiple repetitions throughout training, with model accuracy reaching saturation after only a few passes. Specifically, several aspects remain unexplored:

\begin{itemize}[topsep=0pt]
    \item The influence of arXiv-derived data on targeted mathematical tasks not encompassed by this work, such as the transformation of formal statements into informal prose (informalization of theorems);
    \item The potential interaction effects between arXiv data and other categories of training corpora;
    \item Whether scaling model size may reveal latent advantages of arXiv-based pre-training. The chosen benchmarks represent a range of quantitative reasoning tasks—such as GSM8K \citep{gsm8k}, MATH \citep{MATH}, and CMATH \citep{wei2023cmath}—as well as multiple-choice assessments including MMLU-STEM \citep{mmlu} and Gaokao-MathQA \citep{agieval}. \item \textbf{Direct Preference Optimization (DPO)}: DPO applies additional fine-tuning to the SFT model via augmented output samples, leveraging a pairwise DPO loss for optimization. Consequently, we investigate the approach of iterative RL with GRPO. One plausible explanation is that the DeepSeek-LLM 1.3B model may not possess sufficient capacity to effectively integrate both code and mathematical content in a concurrent training setup. Across the board, mathematical reasoning did not benefit from exposure to arXiv documents; indeed, model performance on a range of mathematical benchmarks either stagnated or declined. Moreover, incorporating \textbf{efficient inference techniques} \citep{xia-etal-2023-speculative,leviathan2023fast,kwon2023efficient,xia2024unlocking}, which dictate how effectively the policy explores, will also be crucial for further progress. Future efforts will include deploying our RL approach on question prompts outside the original data distribution, alongside experimenting with \textbf{advanced sampling (decoding) strategies}, for instance, those informed by tree-search techniques \citep{yao2023tree}. \subsubsection{Evaluation Results}

\textbf{DeepSeekMath~Corpus exhibits exceptional quality, encompasses multilingual mathematical data, and is the largest of its kind.}

\begin{itemize}[topsep=0pt]
    \item \textbf{High-quality}:  
    The model's downstream capabilities are benchmarked across 8 mathematical datasets using few-shot chain-of-thought prompting \cite{cot}. Therefore, our investigation will focus on reinforcement learning techniques that maintain robustness in the presence of noisy or inaccurate reward signals. Incorporating code pretraining leads to improved outcomes on program-assisted mathematical reasoning tasks, observable under both two-stage and one-stage paradigms. The policy model is updated once following each sampling cycle. For benchmark entries containing fewer than 10 but at least 3 grams, we use an exact string match to identify and exclude compromised web pages. This design aims to evaluate whether initializing with code-specific data provides unique benefits for mathematical reasoning relative to a broad general-domain pretraining. Our analysis reveals: 1) With chain-of-thought reasoning, DeepSeekMath-RL 7B achieves 88.2\% and 51.7\% accuracy on GSM8K and MATH, respectively. We omit SFT questions from other domains in order to specifically analyze the influence of RL on benchmarks for which no corresponding training samples are present during the RL stage. According to Figure \ref{fig:maj-pass}, RL provides a boost to Maj@K accuracy, but does not affect Pass@K. For the purposes of this discussion, references to the DeepSeekMath Corpus denote an 89B-token collection assembled in the second round of data curation. To counteract the potential overfitting to the reward model, it is standard practice to incorporate a per-token KL penalty, calculated relative to a reference policy, within the reward signal at each token step \citep{ouyang2022training}, that is,

\begin{equation}
    r_{t} = r_\varphi(q, o_{\le t}) - \beta \log\frac{\pi_{\theta}(o_{t}|q, o_{<t})}{\pi_{ref}(o_{t}|q, o_{<t})},
\label{eq:PPO-reward}
\end{equation}

where $r_\varphi$ denotes the output of the reward model, $\pi_{ref}$ is a reference policy (typically the initial SFT model), and $\beta$ is the weighting for the KL term. The value function in PPO is usually instantiated as an additional model with size comparable to that of the policy model, thereby imposing notable memory and computational costs. Our model's performance is compared against prominent contemporary models:

\begin{itemize}[topsep=0pt]
    \item \textbf{Closed-source models} surveyed include: (1) members of the GPT series, most notably GPT-4 \citep{gpt4} and GPT-4 Code Interpreter~\footnote{\url{https://openai.com/blog/chatgpt-plugins\#\#code-interpreter}}, (2) Gemini Ultra and Pro \citep{gemini}, (3) Inflection-2 \citep{inflection-2}, (4) Grok-1~\footnote{\url{https://x.ai/model-card}},
    as well as recently introduced models from Chinese developers such as (5) Baichuan-3~\footnote{\url{https://www.baichuan-ai.com}},
    and (6) the most recent GLM-4~\footnote{\url{https://open.bigmodel.cn/dev/api\#glm-4}}

    belonging to the GLM family \citep{glm}. \subsubsection{From PPO to GRPO} Proximal Policy Optimization (PPO) \citep{schulman2017proximal} serves as a widely adopted actor-critic RL algorithm during RL-based fine-tuning of LLMs \citep{ouyang2022training}. Furthermore, due to inherent constraints on model size, DeepSeekMath~lags behind GPT-4 with respect to few-shot learning: while GPT-4 demonstrates notable improvement under few-shot conditions, DeepSeekMath~shows little differentiation between its zero-shot and few-shot performance. \item Drawing on this unified perspective, we investigate the mechanisms underpinning the success of reinforcement learning, and outline several promising research directions for achieving more robust reinforcement learning in large language models. To investigate the impact of code training on mathematical reasoning, we conducted both a two-stage and a single-stage training regimen:

\textbf{Two-Stage Training}

\begin{itemize}[topsep=0pt]
    \item \textbf{Code Training for 400B Tokens $\rightarrow$ Math Training for 150B Tokens}: The DeepSeek-LLM 1.3B model is first exposed to 400B tokens of code data, then subsequently trained on 150B math-specific tokens;
    \item \textbf{General Training for 400B Tokens $\rightarrow$ Math Training for 150B Tokens}: For comparative analysis, we replace the initial code data with 400B tokens sampled from a general domain corpus (assembled by DeepSeek-AI) before proceeding with the same 150B math token training. Figure \ref{fig:corpora_comparisons} depicts that DeepSeek-LLM 1.3B, when optimized on DeepSeekMath~Corpus, not only learns more rapidly but also maintains consistent progress over time. \subsection{Training and Evaluating DeepSeekMath-RL}

Our RL experiments are conducted using DeepSeekMath-Instruct 7B as the policy model. To enhance the seed corpus and refine the fastText model, we expand our search to incorporate more mathematical web domains. The fine-tuned DeepSeekMath-Instruct 7B surpasses all other 7B models and even rivals instruction-tuned models with 70B parameters in the open-source community. \paragraph{Mathematical Problem Solving with Step-by-Step Reasoning}

We assess DeepSeekMath-Base on its capability to address mathematical questions via few-shot chain-of-thought prompting \citep{cot} across eight different datasets in both English and Chinese. Following the procedure in \cite{dsp_proof}, we task the model with creating proof outlines, then employ the off-the-shelf automated prover Sledgehammer \citep{sledgehammer} to complete the finer proof details. Allowing models to combine natural language-based reasoning with programmatic tool usage in evaluation further increases DeepSeekMath-Instruct 7B’s accuracy on MATH to nearly 60\%, establishing it as the top-performing open-source model in this setting. The quantity of data preserved is determined by pre-training experiments conducted on the top 40B, 80B, 120B, and 160B tokens, with the first cycle preserving the top 40B tokens. \end{itemize}

\subsubsection{Training Setting}
\label{sec:quality-policy}
We fine-tune a general-purpose language model containing 1.3B parameters for mathematical tasks, utilizing the same architecture as DeepSeek LLMs \citep{deepseek-llm}, hereafter referred to as DeepSeek-LLM 1.3B. The learning rate is scheduled in multiple phases: after a 2,000-step warmup, it ascends to its maximum value, subsequently dropping to 31.6\% of this maximum at 80\% completion of training, and finally tapering off to 10.0\% by the end of 90\% of the training. According to Table \ref{tab:base_math_pot_proof}, DeepSeekMath-Base 7B outstrips the previously best-performing Llemma 34B on these benchmarks. The process rewards are also normalized within the group by their mean and standard deviation, i.e., $\widetilde{r}_i^{index(j)} = \frac{r_i^{index(j)} - {\rm mean(\mathbf{R})}}{{\rm std(\mathbf{R})}}$. Additionally, LLMs have become valuable tools for human users in tackling intricate mathematical challenges~\citep{tao}. The Rule-based function evaluates responses solely based on the correctness of the answer, while the Model-based function employs a learned reward model that assigns scores to each response using training data generated according to the rule-based evaluation. The resulting program output is then used as the answer. \paragraph{Algorithms} Algorithms utilize both the observed data and reward signals to compute gradients, thereby adjusting the model parameters. Finally, we elucidate how our RL approach enhances the performance of instruction-tuned models, and discuss future research avenues for optimizing RL effectiveness within this unified context. DeepSeekMath-RL 7B is evaluated across the same benchmarks as DeepSeekMath-Instruct 7B. \item We establish a unified framework that provides insights into various techniques, including RFT, DPO, PPO, and GRPO. \end{document} \item \textbf{Natural Language Understanding, Reasoning, and Code}: For a thorough examination of the model's linguistic, reasoning, and coding competencies, DeepSeekMath-Base is evaluated against the Massive Multitask Language Understanding (MMLU) benchmark \citep{mmlu}, which includes 57 multiple-choice tasks across various domains, and BIG-Bench Hard (BBH) \citep{bbh}, containing 23 complex tasks requiring multi-step reasoning. Unlike traditional methods, GRPO eliminates the need for a critic model by leveraging group scores to estimate the baseline, which leads to a substantial reduction in the computational resources necessary for training compared to Proximal Policy Optimization (PPO). Furthermore, within reinforcement learning training, this value function acts as a baseline in advantage estimation to reduce variance. Specifically, preliminary testing revealed challenges in handling questions involving triangles and ellipses, suggesting a possible bias introduced during both data pre-selection and downstream fine-tuning. DeepSeekMath-Base leverages DeepSeek-Coder-Base-v1.5 7B \citep{deepseek-coder} as its initialization point, as our analysis reveals that models pre-trained on code exhibit superior performance relative to those based on general-purpose LLMs. The model generates formal proofs in Isabelle for each instance using few-shot prompting. Policy optimization is again performed

\subsubsection{Iterative RL with GRPO} During the progression of reinforcement learning, the previous reward model may no longer adequately guide the latest policy model. \subsubsection{Process Supervision RL with GRPO} 
While outcome supervision delivers a reward only at sequence completion, this granularity may fall short for supervising agents in intricate mathematical problems. Through a comprehensive ablation analysis, we demonstrate that web-derived data represent a rich reservoir for high-quality mathematical resources, whereas arXiv contributions appear less impactful than initially anticipated. \begin{itemize}[topsep=0pt]
    \item \textbf{English mathematical datasets}: We provide tool-integrated solution annotations for problems from GSM8K and MATH, and incorporate a selected portion of MathInstruct \citep{MathInstruct} in addition to the training set from Lila-OOD \citep{lila}, wherein solutions are given using either CoT or PoT strategies. \end{itemize}

A summary of the method components is presented in Table \ref{tab:data_policy}. This advantage holds even relative to much larger systems (for example, Qwen 72B) or those fine-tuned for mathematical reasoning using reinforcement learning (such as WizardMath-v1.1 7B). To curtail data volume, URL-based deduplication and near-duplicate removal are applied to Common Crawl, reducing the dataset to 40B HTML web pages. Our process starts with OpenWebMath \citep{openwebmath}, a carefully curated compilation of mathematical web content, as the initial seed dataset. Training employs a batch size of 4M tokens and supports sequences of up to 4K tokens in context length. \subsection{Reflections on Reinforcement Learning}
\subsubsection{Toward a Unified Framework} In what follows, we introduce a cohesive framework to compare and analyze various training approaches—including SFT, RFT, DPO, PPO, and GRPO—and design experiments to investigate key determinants within this paradigm. We evaluate DeepSeekMath-Base 7B on the informal-to-formal proof generation task introduced by \citep{dsp_proof}, where the objective is to synthesize a formal proof from an informal assertion, its corresponding formalized version, and an informal reasoning chain. \section{Math Pre-Training}

\subsection{Data Collection and Decontamination} This section describes the methodology used to assemble the DeepSeekMath Corpus from Common Crawl. Within our experiments, we use two variants for reward functions: `Rule' and `Model'. This English dataset collection spans an array of mathematical disciplines, such as algebra, probability, number theory, calculus, and geometry. At the beginning of training, Online RFT exhibits results similar to RFT, but as training progresses, it increasingly surpasses RFT, highlighting the benefits of an online training regime. \subsection{Contributions}
This work presents advancements in large-scale math pre-training, alongside comprehensive studies in reinforcement learning. We further investigate iterative RL by performing two rounds of training; Figure \ref{fig:iter-rl} reveals that iterative RL produces marked improvements, with the first iteration yielding particularly significant gains. Our filtration procedure entails removing any text where a 10-gram phrase precisely matches a substring from the evaluation benchmarks, eliminating these segments from our math training data. Moreover, we plan to investigate the research paths highlighted in Section \ref{sec:effec_rl} for further advancing reinforcement learning methodologies in LLMs. Following the initial round of data collection, a considerable number of mathematical web pages are still absent, primarily due to the limited diversity within the positive training examples for the fastText model. \section{Conclusion, Limitation, and Future Work} 

In this work, we introduce DeepSeekMath, a model that establishes new benchmarks for open-source systems on competition-level MATH tasks and attains near-parity with leading proprietary models. \paragraph{Formal Mathematics}

Recent years have seen increasing emphasis on formal proof automation, which enhances both the accuracy and efficiency of mathematical verification. Additionally, we expand the scope of evaluation to general abilities, assessing the model’s competency in natural language processing, logical reasoning, and software development. This pattern echoes the findings of \citep{wang2023making}, who also reported a \textbf{misalignment problem} for reasoning tasks within SFT models and demonstrated that various preference alignment techniques can successfully improve the reasoning performance of these models \citep{yuan2023rrhf,song2023preference,wang2023making}. In addition, we note improvements in out-of-domain generalization during reinforcement learning. \item \textbf{Rejection Sampling Fine-tuning (RFT)}: Building on the SFT model, RFT continues fine-tuning using SFT-derived questions, while filtering model outputs to retain only those with correct responses. We propose three critical research directions concerning reward models: 1) \textbf{Enhancing the generalization capabilities of the reward model.} For reinforcement learning to produce fundamental advances—rather than simply reinforce the status quo of large language model (LLM) outputs—the reward model should demonstrate strong generalization to out-of-distribution queries and sophisticated generations; 2) \textbf{Capturing and representing reward model uncertainty.} Quantifying uncertainty may provide a crucial interface linking imperfect reward signals to weak-to-strong alignment strategies; 3) \textbf{Developing efficient approaches to construct high-quality process reward models} that yield granular and informative signals for the reasoning steps during model training \citep{lightman2023let,wang2023math}. Evaluation is conducted using miniF2F \citep{minif2f}, a dataset consisting of formal mathematics problems at Olympiad difficulty. Nevertheless, the reliability of reward signals cannot be guaranteed at all times, particularly in tasks characterized by substantial complexity. In RL, we particularly define the data source as consisting of unlabeled problems paired with sampled responses generated by the policy model. We outline possible research avenues related to each component. This suggests that model size is not the sole determinant of mathematical reasoning capability; high-quality pre-training on domain-specific data can enable smaller models to perform competitively. Across other datasets, its results are competitive with DeepSeek-LLM-Chat 67B, which previously represented the open-source state-of-the-art at an order of magnitude larger model size. DeepSeekMath-Base achieves robust few-shot results in automatic formalization. \subsubsection{ArXiv Papers Appear Not to Enhance Mathematical Reasoning}

Although arXiv papers are widely utilized as a component within math pre-training datasets \citep{minerva,gpt-f,llemma,mathpile}, there remains little granular analysis concerning their specific contribution to mathematical reasoning skills. For DeepSeekMath-RL 7B, the GSM8K and MATH datasets featuring chain-of-thought solutions are treated as in-domain, while all remaining evaluation sets are categorized as out-of-domain. According to Equation \ref{eq:objective}, existing approaches generally \textbf{TRUST} the reward signal, using it to either encourage or discourage the model's likelihood of generating specific tokens. Experimental evidence indicates GRPO continues to offer meaningful gains, even when DeepSeekMath-Instruct 7B performs strongly on standard benchmarks. The English testbeds comprise GSM8K \citep{gsm8k}, MATH \citep{MATH}, SAT \citep{llemma}, OCW Courses \citep{minerva}, and MMLU-STEM \citep{mmlu}. \subsubsection{How to Achieve More Effective RL?}
\label{sec:effec_rl}
Our results confirm that RL is highly effective for mathematical reasoning. The data composition for training includes 56\% from the DeepSeekMath Corpus, 4\% sourced from AlgebraicStack, 10\% from arXiv, 20\% comprising Github code, with the final 10\% being natural language content drawn from the English and Chinese portions of Common Crawl. The assessment criteria include the models’ proficiency in generating independent textual solutions (without recourse to external tools) and their capability to tackle problems programmatically via Python. For each prompt, $64$ candidate completions are generated. \section{Discussion}
This section details our insights gained from pre-training and RL experiments. The classifier is then utilized to identify further positive examples within CC, which are subsequently reviewed and refined by human annotators. The parameter $\varepsilon$ serves as a clipping hyper-parameter in PPO, aiding in the stabilization of the training process. Figure \ref{fig:corpora_comparisons} further illustrates that, after training on 50B tokens (representing a full epoch for Proof-Pile-2), models using the DeepSeekMath Corpus still surpass those trained on Proof-Pile-2, suggesting an overall superior average quality of DeepSeekMath Corpus. In addition, we present Group Relative Policy Optimization (GRPO), an adaptation of the Proximal Policy Optimization (PPO) reinforcement learning algorithm \citep{schulman2017proximal}. Those domains in which more than 10\% of pages have been collected are marked as being math-oriented (such as \url{mathoverflow.net}). On the English benchmarks, DeepSeekMath-Base achieves performance that rivals the proprietary Minerva 540B \citep{minerva} model, and it notably outperforms all existing open-source base models such as Mistral 7B \citep{mistral} and Llemma-34B \citep{llemma}, regardless of their exposure to math-specific pre-training—frequently by a notable margin. Each problem is accompanied by its solution formatted as chain-of-thought (CoT) \citep{cot}, program-of-thought (PoT) \citep{pot,pal}, or tool-integrated reasoning explanations \citep{tora}. Subsequent to pre-training, we carry out mathematical instruction tuning on DeepSeekMath-Base using datasets featuring chain-of-thought \citep{cot}, program-of-thought \citep{pot,pal}, and tool-integrated reasoning \citep{tora} strategies. \item \textbf{Multilingual}:  
    The DeepSeekMath~Corpus is characterized by its diverse linguistic content, with English and Chinese comprising the largest portions. In this section, we introduce Group Relative Policy Optimization (GRPO), an RL algorithm we developed for efficiency and effectiveness. \item \textbf{PPO/GRPO}: PPO/GRPO initializes its policy with the SFT model, then applies reinforcement learning by sampling outputs from the policy as it updates in real time. Next, the fastText model retrieves math-relevant pages from this deduplicated corpus. \end{itemize}

\subsection{Training and Evaluating DeepSeekMath-Instruct 7B}

This section presents the development and evaluation process of DeepSeekMath-Instruct 7B, which is instruction-tuned for mathematics tasks using DeepSeekMath-Base as its foundation. By the fourth iteration, we find that approximately 98\% of all available data was already captured in the third iteration, prompting us to terminate further collection. The development of DeepSeekMath is underpinned by the construction of the DeepSeekMath Corpus—a massive, high-quality pre-training collection containing 120B math-focused tokens. \end{itemize}

\paragraph{Results}
Tables \ref{tab:code-math} and \ref{tab:code-math-general-eval} report downstream outcomes for the various training protocols. By utilizing a group-based comparative framework for advantage estimation, GRPO complements the comparative structure on which reward models are commonly trained (i.e., comparisons across outputs to the same prompt). Nonetheless, top-performing systems like GPT-4~\citep{gpt4} and Gemini-Ultra~\citep{gemini} remain proprietary, and publicly released open-source counterparts still exhibit a substantial gap in performance. \end{itemize}

\textbf{Exploration and Analysis of Reinforcement Learning}

\begin{itemize}[topsep=0pt]
    \item We propose Group Relative Policy Optimization (GRPO), a novel reinforcement learning approach characterized by both high efficiency and effectiveness. Table \ref{tab:corpora_comparison} highlights that models trained on the DeepSeekMath~Corpus consistently outperform those trained on alternative corpora. It is important to emphasize that the pipeline is extensible to additional domains, including programming. Online sampling refers to data collected from the real-time exploration of the current policy model during training. To minimize benchmark leakage, we adopt the strategy from \cite{deepseek-coder}, filtering out any web page containing content from English mathematical benchmarks, such as GSM8K \citep{gsm8k} and MATH \citep{MATH}, as well as from Chinese datasets such as CMATH \citep{wei2023cmath} and AGIEval \citep{agieval}. After conducting four iterative data collection rounds, we obtain 35.5M mathematical web pages, totaling 120B tokens. The RL training corpus comprises approximately 144K chain-of-thought questions drawn from GSM8K and MATH datasets present in the SFT data. Our comprehensive experiments—including comparisons of online versus offline training, outcome versus process-based supervision, and single-turn versus iterative RL—enable an in-depth exploration of core components underlying this paradigm. Offline sampling, in contrast, utilizes data generated by the preliminary SFT model prior to further training. We also systematically evaluate factors such as online versus offline optimization, supervision at the outcome versus process level, single-step versus iterative reinforcement learning, among others, to thoroughly examine the foundational aspects of this paradigm. Referencing Table \ref{tab:sft_rl_math}, when evaluated without access to external tools, DeepSeekMath-Instruct 7B excels at stepwise mathematical reasoning. As the training continues, data sampled from the evolving actor diverges more substantially, amplifying the advantages of online data collection. These results suggest that RL contributes to enhancing the robustness of the output distribution, specifically by increasing the likelihood that the correct response appears among the TopK candidates, rather than directly upgrading the model’s inherent capabilities. Given the prompt $q$ and $G$ samples $\{o_1, o_2, \cdots, o_G\}$, a process-based reward model evaluates each individual step, producing structured rewards: $\mathbf{R} = \{ \{r_1^{index(1)},\cdots,r_1^{index(K_1)}\}, \cdots,  \{r_G^{index(1)},\cdots,r_G^{index(K_G)}\} \}$, where $index(j)$ is the index of the last token in step $j$, and $K_i$ is the total step count for $o_i$. DeepSeekMath~7B demonstrates a notable achievement by attaining a 51.7\% score on the MATH benchmark—designed for competition-level assessment—without the aid of supplementary toolkits or ensemble approaches, reaching performance on par with Gemini-Ultra and GPT-4. This scalar value is then uniformly assigned as the advantage $\hat{A}_{i, t}$ for every token in $o_i$, namely, $\hat{A}_{i, t} = \widetilde{r}_i = \frac{r_i- {\rm mean}(\mathbf{r})}{{\rm std}(\mathbf{r})}$. \subsubsection{Why RL Works?}  
In this study, we utilize reinforcement learning on a subset of the instruction tuning dataset, and observe considerable improvements over the baseline instruction tuning model. Table \ref{tab:sft_rl_math} presents a comparison of open- and closed-source models, examining their performance on English and Chinese benchmarks using both chain-of-thought and tool-integrated reasoning. Chinese evaluation suites include MGSM-zh \citep{mgsm}, CMATH \citep{wei2023cmath}, Gaokao-MathCloze \citep{agieval}, and Gaokao-MathQA \citep{agieval}. In sum, across all three benchmarks in reasoning and code generation, DeepSeekMath-Base 7B substantially outperforms the generalist Mistral 7B model \citep{mistral}. This work presents DeepSeekMath, a specialized language model tailored for mathematics that not only surpasses existing open-source alternatives in mathematical reasoning but also nears GPT-4's proficiency on academic evaluation tasks. This work presents DeepSeekMath~7B, an extension of DeepSeek-Coder-Base-v1.5 7B, further pre-trained on a substantial 120B tokens consisting of mathematical content from Common Crawl, along with related natural language and code corpora. For a comprehensive explanation of the derivation steps, please consult Appendix \ref{app:analysis-rl}. Furthermore, by including code tokens during continued training, DeepSeekMath-Base 7B preserves the coding performance of DeepSeek-Coder-Base-v1.5 on HumanEval and MBPP. \item While pre-training on arXiv documents is prevalent in math-related research, we observe no significant performance gain across all mathematical benchmarks included in this study when incorporating such data. Perhaps unexpectedly, our experimental findings indicate that arXiv papers do not appear to facilitate improved mathematical reasoning capabilities. DeepSeekMath~is derived from DeepSeek-Coder-v1.5 7B through extended training over 500B tokens, including a substantial 120B token subset of mathematical content curated from Common Crawl.